% INLINED into encounter_multimodal_prr_v2.tex on 2026-09-03 (Sec. III); this file is no
% longer \input anywhere and is kept for reference only -- edit the master file instead.

\section{Prescribed finite modality: model and theorem}
\label{sec:exact-m-spine}

The conserved-budget support construction admits encounter processes with any
prescribed \emph{fixed finite} number of reaction-time modes.  We state this
existence result in an exact quotient of a two-particle Doi process; the
construction depends on the prescribed dimension and mode count and does not
assert mode-count switching by weight redistribution on one fixed support.

Fix finite $d\ge2$, finite $m\ge1$, a transverse period $W>0$, a contact
radius $0<a<W/2$, positive scales $\ell_0,\rho>0$, and a compact interval
$I=[\tau,T]\subset(0,\infty)$.
On $\mathbb R\times\mathbb R\times\mathbb T_W^{d-1}$, consider the
independent midpoint and relative processes
\begin{equation}
 \begin{aligned}
 \dd Z_t&=-\gamma(Z_t-\bar z)\,\dd t
       +\varepsilon\sqrt{D_0}\,\dd W_t,\\
 \dd R_{\parallel,t}&=-\gamma R_{\parallel,t}\,\dd t
       +2\varepsilon\sqrt{D_0}\,\dd W_{\parallel,t},\\
 \dd R_{\perp,t}&=2\varepsilon\sqrt{D_0}\,\dd W_{\perp,t}\pmod W.
 \end{aligned}
 \label{eq:exact-m-process}
\end{equation}

The mutually independent initial laws satisfy
\begin{equation}
 \begin{aligned}
 Z_0&\sim N\!\left(z_0,\varepsilon^2D_0/(2\gamma)\right),\\
 R_{\parallel,0}&\sim N(r_{\parallel,0},\varepsilon^2u_0^2),\\
 R_{\perp,0}&\sim
 \operatorname{WN}(r_{\perp,0},\varepsilon^2\Sigma_{\perp,0}),
 \end{aligned}
 \qquad
 \begin{gathered}
 D_0,\gamma,u_0>0,\quad z_0\ne\bar z,\\
 u_0^2<4D_0/\gamma,\quad \Sigma_{\perp,0}\succ0.
 \end{gathered}
 \label{eq:exact-m-initial}
\end{equation}

Here $\operatorname{WN}$ is the wrapped Gaussian on
$\mathbb T_W^{d-1}$, $\Sigma_{\perp,0}$ is a positive-definite
$(d-1)\times(d-1)$ covariance matrix, and $|r|_{\rm mi}$ is the Euclidean norm
of the longitudinal relative coordinate together with the shortest transverse
representatives modulo $W$.

All driving Brownian motions are also independent.  Thus
$\mu(t)=\bar z+(z_0-\bar z)\e^{-\gamma t}$ is strictly monotone and
$\operatorname{Var}Z_t=\varepsilon^2D_0/(2\gamma)$: the midpoint is prepared
with its stationary variance so that only its mean relaxes, which makes every
exposure clock below have the same width $\sigma$.  The inequality
$u_0^2<4D_0/\gamma$ is used only to place the initial law in the weighted
state space of the transfer estimate in the Supplemental Material.  Under the
minimum-image
convention, assume the deterministic relative trajectory
$r_*(t)=(r_{\parallel,0}\e^{-\gamma t},r_{\perp,0})$ obeys
$\sup_{t\in I}|r_*(t)|_{\rm mi}\le a-\eta_*$ for some $\eta_*>0$.

Choose target times
\begin{equation}
 \begin{aligned}
  &\tau<t_1<\cdots<t_m<T,\qquad c_j=x(t_j),\\
  &x(t)=\frac{\operatorname{sgn}(\mu')\mu(t)}{\ell_0},
 \end{aligned}
 \label{eq:exact-m-centres}
\end{equation}
so that $x'(t)$ is bounded away from zero.  At the corresponding physical
centres, place normalized Gaussian catalyst slabs of width
$\varepsilon\rho$.  Their nonnegative weights belong to a compact subset of
the simplex,
\begin{equation}
 \mathcal W_{w_*}
 =\left\{\bm w:\
   \substack{\sum_{j=1}^{m}w_j=1,\\ w_j\ge w_*>0}\right\},
 \qquad 0<w_*\le\frac1m,
 \label{eq:exact-m-simplex}
\end{equation}
and define
\begin{equation}
 \begin{aligned}
 \phi_{j,\varepsilon}(z)
 &=\frac{\exp[-(z-\mu(t_j))^2/(2\varepsilon^2\rho^2)]}
        {\sqrt{2\pi}\,\varepsilon\rho},\\
 K_{B,\bm w,\varepsilon}(Z,R)
 &=\bm1_{\{|R|_{\rm mi}<a\}}\frac{B}{W^{d-1}}
   \sum_{j=1}^{m}w_j\phi_{j,\varepsilon}(Z).
 \end{aligned}
 \label{eq:exact-m-killing}
\end{equation}

The slab normalization makes $B$ the same installed centre-space budget for
every allocation.  Set $K_{B,\bm w,\varepsilon}=B V_{\bm w,\varepsilon}$,
and let $T_0(t)$ and $q_0$ denote the unkilled semigroup and initial law.  Then
\begin{equation}
 \begin{aligned}
 G_{\varepsilon,\bm w}(t)
 &=B^{-1}\mathbb E_0[K_{B,\bm w,\varepsilon}(Z_t,R_t)]
 =\langle V_{\bm w,\varepsilon},T_0(t)q_0\rangle,\\
 f_{B,\varepsilon,\bm w}(t)
 &=\mathbb E_0\!\left[K_{B,\bm w,\varepsilon}(Z_t,R_t)
  \e^{-\int_0^tK_{B,\bm w,\varepsilon}(Z_s,R_s)\,\dd s}\right].
 \end{aligned}
 \label{eq:exact-m-observables}
\end{equation}

The expectation $\mathbb E_0$ is with respect to the unkilled process started
from Eq.~\eqref{eq:exact-m-initial}.  Independence gives the exact
factorization
\begin{align}
 G_{\varepsilon,\bm w}(t)
 &=a_\varepsilon(t)H_{\sigma,\bm w}(x(t)),\nonumber\\
 a_\varepsilon(t)
 &=\frac{c_{d,\varepsilon}(t)}
 {W^{d-1}\sqrt{2\pi}\,\varepsilon
  \sqrt{D_0/(2\gamma)+\rho^2}},
 \label{eq:exact-m-factorization}
\end{align}
where $c_{d,\varepsilon}(t)=
\mathbb P_0(|R_t|_{\rm mi}<a)$ and $H_{\sigma,\bm w}$ is defined below.
All dependence on $d$ enters through $c_{d,\varepsilon}$ and the
normalization $W^{-(d-1)}$.

\medskip
\noindent\textit{Theorem 1 (exact prescribed finite modality).---}
Fix finite $d\ge2$, finite $m\ge1$, the data above, and the compact window
$I$.
(i) There is $\varepsilon_0>0$ such that, for every
$0<\varepsilon<\varepsilon_0$ and every $\bm w\in\mathcal W_{w_*}$, the
unit-budget free-exposure clock $G_{\varepsilon,\bm w}$ has exactly $m$
nondegenerate interior maxima and $m-1$ nondegenerate interior minima on
$I$, ordered alternately, and no stationary endpoint.
(ii) For each such $\varepsilon$ there is a budget
$B_{\rm top}^{\rm unif}(\varepsilon)>0$ such that the budget-rescaled Doi
reaction-time density $f_{B,\varepsilon,\bm w}/B$ has the same complete
stationary signature on $I$ for every
$0<B<B_{\rm top}^{\rm unif}(\varepsilon)$ and every
$\bm w\in\mathcal W_{w_*}$.  Multiplication by $B$ leaves the stationary
points unchanged, so the conclusion also holds for the physical density
$f_{B,\varepsilon,\bm w}$.

Precisely, for each fixed such $\varepsilon$ and allocation set
$\mathcal A\subseteq\mathcal W_{w_*}$, define the topology-retention budget
\begin{align}
 B_{\rm top}(\varepsilon;\mathcal A)
 &:=\sup\bigl\{b>0:\nonumber\\[-2pt]
 &\quad f_{\beta,\varepsilon,\bm w}/\beta
 \text{ has the stated}\nonumber\\[-2pt]
 &\quad\text{complete stationary signature on }I
 \nonumber\\[-2pt]
 &\quad\text{for every }\bm w\in\mathcal A
 \text{ and every }0<\beta<b\bigr\}.
 \label{eq:exact-m-btop}
\end{align}
Then
\begin{equation}
 B_{\rm top}^{\rm unif}(\varepsilon)
 :=B_{\rm top}(\varepsilon;\mathcal W_{w_*})>0.
 \label{eq:exact-m-btop-unif}
\end{equation}
and (ii) holds with this value.

The mechanism can be seen without invoking a normal form.  Gaussian
convolution reduces the longitudinal part of the unit-budget exposure to
\begin{equation}
 \begin{aligned}
  H_{\sigma,\bm w}(x)
  &=\sum_{j=1}^{m}w_j
    \exp\!\left[-\frac{(x-c_j)^2}{2\sigma^2}\right],\\
  \sigma&=\frac{\varepsilon}{\ell_0}
  \sqrt{\frac{D_0}{2\gamma}+\rho^2}.
 \end{aligned}
 \label{eq:exact-m-mixture}
\end{equation}

Writing $q_j(x)=w_j\exp[-(x-c_j)^2/(2\sigma^2)]$,
$\pi_j(x)=q_j(x)/H_{\sigma,\bm w}(x)$, and
$L=\partial_x\log H$, let $\bar c(x)$ and
$\operatorname{Var}_{\pi(x)}(c)$ denote the mean and variance of the centres
under the finite probability vector $\bm\pi(x)$.  Then
\begin{equation}
 L(x)=\frac{\bar c(x)-x}{\sigma^2},
 \qquad
 L'(x)=\frac{\operatorname{Var}_{\pi(x)}(c)}{\sigma^4}
       -\frac{1}{\sigma^2}.
 \label{eq:exact-m-log-slope}
\end{equation}

The univariate upper bound of $m$ modes for an $m$-component Gaussian mixture
is classical \cite{carreiraPerpinanWilliams2003modes,
amendolaEngstromHaase2020modes}.  Here a direct affine--exponential proof is
retained because the Doi transfer needs all critical zeros counted with
multiplicity together with uniform root-tube, complement, curvature, and
endpoint margins.

For the adjacent pair $c_j<c_{j+1}$, the weighted crossover is
\begin{equation}
 s_j=\frac{c_j+c_{j+1}}{2}
 +\frac{\sigma^2}{c_{j+1}-c_j}
  \log\!\frac{w_j}{w_{j+1}}.
 \label{eq:exact-m-crossover}
\end{equation}

The adjacent ratio obeys
$q_{j+1}/q_j=1/9$ at
$s_j-\sigma^2\log(9)/(c_{j+1}-c_j)$ and $q_{j+1}/q_j=9$ at the corresponding
plus edge.  Thus both adjacent components remain uniformly present throughout
the crossover, while all nonadjacent components are exponentially small.
Equation~\eqref{eq:exact-m-log-slope} then yields one simple minimum at
\begin{equation}
 r_j=s_j+O(\sigma^4),
 \qquad
 L'(r_j)=\frac{(c_{j+1}-c_j)^2}{4\sigma^4}
          +O(\sigma^{-2}).
 \label{eq:exact-m-valley}
\end{equation}

Each centre similarly contributes one simple maximum at
$p_j=c_j+O(\exp[-q/\sigma^2])$, with
$L'(p_j)=-\sigma^{-2}+o(\sigma^{-2})$.  These constructions give
$2m-1$ simple stationary points.  They are exhaustive: after multiplication
by a positive Gaussian factor, $H'$ is an exponential polynomial
\begin{equation}
 \sum_{j=1}^{m}(a_j+b_jx)\e^{\lambda_jx},
 \qquad \lambda_1<\cdots<\lambda_m,
 \label{eq:exact-m-exp-polynomial}
\end{equation}
which has at most $2m-1$ real zeros counted with multiplicity.  The extreme
affine--exponential terms dominate in the two tails, so all zeros lie in a
compact interval and real analyticity makes their number finite.  A
generalized Rolle induction then proves the bound by differentiating twice to
remove the lowest-exponent affine term.

The remaining relative-coordinate contact probability is a positive slow
factor.  Uniform bounds on its first two logarithmic time derivatives, peak
boxes of width $O(\sigma^2)$, valley boxes of width $O(\sigma^4)$, and a
sign argument on the complement of these boxes (the posterior-sector
certificate of the Supplemental Material: a sign bound on the logarithmic
slope $L$ obtained from the relative weights $\pi_j(x)$ of the slabs)
exclude any further stationary point.  Consequently the slow factor moves maxima by $O(\sigma^2)$ and
minima by $O(\sigma^4)$ without changing their number or type.  Finally, on
every fixed positive $\varepsilon$, convergence of the budget-rescaled
density together with its first two time derivatives, uniformly in the slab
weights (Sec.~\ref{sec:transfer}),
\begin{equation}
 \sup_{\bm w\in\mathcal W_{w_*}}
 \left\|\frac{f_{B,\varepsilon,\bm w}}{B}
       -G_{\varepsilon,\bm w}\right\|_{C^2(I)}
 \longrightarrow0
 \qquad (B\downarrow0)
 \label{eq:exact-m-weak-budget}
\end{equation}
preserves the disjoint neighbourhoods of the stationary points (root
tubes), their signed curvature margins, the derivative margin on their
complement, and the endpoint slopes.

The quantifier order is essential: the geometry and admissible scales may
depend on fixed $(d,m)$; one first fixes a sufficiently small
$\varepsilon>0$ and only then chooses
$B<B_{\rm top}^{\rm unif}(\varepsilon)$.  The theorem does not provide a
useful lower bound for that threshold, a uniform-in-$d$ or uniform-in-$m$
construction, an event-mass floor, or a stationary-point count outside $I$;
it also uses a whole-window contact-interior condition whose small-noise limit
is asymptotically saturated contact.  Nontrivial contact, observability at a
common positive budget, and finite-parameter robustness remain separate
numerical checks.
